\documentclass{article}
\usepackage{graphicx} % Required for inserting images

\title{Sutton and Barton Exercises Chapter 7}
\author{danieltocco0 }
\date{November 2025}

\begin{document}

\maketitle

\section{Chapter 7}
\\\\
\textit{Exercise 7.1} In Chapter 6 we noted that the Monte Carlo error can be written as the
sum of TD errors (6.6) if the value estimates don’t change from step to step. Show that
the n-step error used in (7.2) can also be written as a sum TD errors (again if the value
estimates don’t change) generalizing the earlier result. 
\\
\\
$G_{t:t+n} - V_{t+n-1}(S_{t}) = R_{t+1} + \gamma G_{t+1:t+n} - V_{t+n-1}(S_{t}) + \gamma V_{t+n-1}(S_{t+1}) - \gamma V_{t+n-1}(S_{t+1}) \\

$= \delta + \gamma (G_{t+1:t+n} - \gamma V_{t+n-1}(S_{t+1})) \Rightarrow$ \\

$\delta + \gamma (G_{t+2:t+n} - \gamma V_{t+n-1}(S_{t+2})) \Rightarrow$ \\

$\sum_{k=t}^{t+n-1}\gamma^{k-t}\delta_{k}$
\\\\

\textit{Exercise 7.3} Why do you think a larger random walk task (19 states instead of 5) was
used in the examples of this chapter? Would a smaller walk have shifted the advantage
to a different value of n? How about the change in left-side outcome from 0 to -1 made
in the larger walk? Do you think that made any difference in the best value of n?  
\\
\\

Because with only 5 states, there are so few states, it is hard to see the 

difference (and benefit) of using n-step methods. The bias-variance tradeoff 

is more pronounced with 19 states. To be sure, there is higher bias in fewer 

step methods and more variance in higher step methods. 
\\

A smaller walk would have shifted the advantage to a smaller n. In particular

for this problem, n=2 would be best. If n=3, then the value of C could be 

pulled towards 1, which would be inaccurate.
\\

No, the optimal value of n would not change, as: "Scaling rewards scales 

the variance up or down, but it doesn't change which n trades off bias vs 

variance best." In other words, the bias-variance structure is maintained.

If one wanted to change the optimal value of n, one could: change the

transition dynamics, add stochastic resets, change the environment to

make early termination more common, or add stochastic reset

(randomly get send back to the starting position).
\\\\

\textit{Exercise 7.4} Prove that the n-step return of Sarsa (7.4) can be written exactly in terms
of a novel TD error, as (too lazy to write it out) 
\\
\\

$G_{t:t+n} - Q_{t-1}(S_{t}, A_{t}) = R_{t+1} + \gamma G_{t+1:t+n} - Q_{t-1}(S_{t}, A_{t}) \Rightarrow$
\\

$G_{t:t+n} - Q_{t-1}(S_{t}, A_{t}) = R_{t+1} + \gamma G_{t+1:t+n} - Q_{t-1}(S_{t}, A_{t}) - \gamma Q_{t}(S_{t+1}, A_{t+1}) + 

\gamma Q_{t}(S_{t+1}, A_{t+1}) 
\\

= \delta_{t} + \gamma (G_{t+1:t+n} - Q_{t}(S_{t+1}, A_{t+1})) \Rightarrow
\\

\delta_{t} + \gamma \delta_{t+1} + \gamma^{2}(G_{t+2:t+n} - Q_{t+1}(S_{t+2}, A_{t+2})) \Rightarrow
\\

\sum_{k=t}^{t+n-1}\gamma^{k-t}\delta_{k} + \gamma^{n}(G_{t+n:t+n} - Q_{t+n-1}(S_{t+n}, A_{t+n}))$
\\

$G_{t+n:t+n} - Q_{t+n-1}(S_{t+n}, A_{t+n})$ is 0, however. We can add $Q_{t-1}(S_{t}, A_{t})$ over

and recover 6.6.
\\\\
\textit{Exercise 7.5} Write the pseudocode for the off-policy state-value prediction algorithm
described above. 
\\
\\

We have to initialize V(s) arbitrarily because we're working with state 

values. 
\\

Next: If $\tau + n < T$, then: $\rho_{\tau}(R_{t+1} + \gamma G_{\tau+1:h}) + (1 - \rho_{\tau}(V_{\tau + n -1}(S_{t}))$ \\

Then: $V(S_{\tau}) \leftarrow V(S_{\tau}) + \alpha[G - V(S_{\tau})]$
\\\\
\textit{Exercise 7.5} Prove that the control variate in the above equations does not change the
expected value of the return.  
\\
\\

The control variate part is: \gamma$Q_{h-1}(S_{t+1}, A_{t+1}) + \gamma \bar{V}_{h-1}(S_{t+1})$
\\

Taking \gamma$\mathbb{E}[Q_{h-1}(S_{t+1}, A_{t+1}) + \bar{V}_{h-1}(S_{t+1})]$ 
\\

= $\sum_{a}\mu(A_{t+1}|S_{t+1})\frac{\pi(A_{t+1}|S_{t+1})}{\mu(A_{t+1}|S_{t+1})}[Q_{h-1}(S_{t+1}, A_{t+1}) + \bar{V}_{h-1}(S_{t+1})]$
\\

= $\sum_{a}\pi(A_{t+1}|S_{t+1})[Q_{h-1}(S_{t+1}, A_{t+1}) + \bar{V}_{h-1}(S_{t+1})]$
\\

Because $\gamma \bar{V}_{h-1}(S_{t+1})$ does is not a probability distribution over actions, 

$\sum_{a}\pi(A_{t+1}|S_{t+1}) \bar{V}_{h-1}(S_{t+1})] = \bar{V}_{h-1}(S_{t+1})$. And now:
\\

$\sum_{a}\pi(A_{t+1}|S_{t+1})[Q_{h-1}(S_{t+1}, A_{t+1})] = \bar{V}_{h-1}(S_{t+1})]$ and so:
\\

$\bar{V}_{h-1}(S_{t+1}) -  \bar{V}_{h-1}(S_{t+1}) = 0.$
\\\\
\textit{Exercise 7.7} Show that the general (off-policy) version of the n-step return (7.13) can
still be written exactly and compactly as the sum of state-based TD errors (6.5) if the
approximate state value function does not change.  
\\
\\

Define the TD error: 

$\delta_{t} = \rho_{t}(R_{t+1} + \gamma G_{t+1:h} - V(S_{t}))$ 
\\

Subtract $V(S_{t})$ so we can form the TD error: 

$G_{t:h} - V(S_{t})= \rho_{t}(R_{t+1} + \gamma G_{t+1:h}) + (1-\rho_{t})V_{h-1}(S_{t}) - V(S_{t})$
\\

$= \rho_{t}(R_{t+1} + \gamma G_{t+1:h} - V(S_{t}) + (1-\rho_{t})(V(S_{t}) - V_{h-1}(S_{t}))$
\\

where above we used the fact that $V(S_{t}) = V(S_{t})(1-\rho_{t}) + \rho V(S_{t})$. 

Now under the assumption of the approximate state value function not 

changing, $V(S_{t}) = V_{h-1}(S_{t})$, we get: 
\\

$G_{t:h} - V(S_{t}) =  \rho_{t}(R_{t+1} + \gamma G_{t+1:h} - V(S_{t}))$
\\

Now $G_{t+1:h}$ can be unrolled as:
\\

$G_{t+1:h} = \rho_{t+1}(R_{t+2} + \gamma G_{t+2:h}) + (1 - \rho_{t+1})V(S_{t+1})$
\\

which we can substitute back into $G_{t:h} - V(S_{t})$:
\\

$\rho_{t}(R_{t+1} + \gamma (\rho_{t+1}(R_{t+2} + \gamma G_{t+2:h}) + (1 - \rho_{t+1})V(S_{t+1}) - V(S_{t}))$. 
\\

For emphasis, note again that V does not change, so $(1-p_{t+1})V(S_{t+1})$ again 

gets cancelled out, and we get:
\\

$\rho_{t}(R_{t+1} + \gamma (\rho_{t+1}(R_{t+2} + \gamma G_{t+2:h}) - V(S_{t}))$. 
\\

This repeats for all intermediate terms, and the last term $V(S_h)$ also 

disappears, as we'd have $G_{h:h} - V(S_{h})$ (with $\gamma$ and $\rho$ factors shared between 

them). We thus ultimately end up with: 
\\

$G_{t:h} - V(S_{t}) = \sum_{t}^{t-1}\gamma^{h-t}(\prod_{k}^{h-1}\rho_{k})\delta_{t}$

\\\\
\textit{Exercise 7.11} Show that if the approximate action values are unchanging, then the
tree-backup return (7.16) can be written as a sum of expectation-based TD errors:
$G_{t:t+n} = Q(S_{t}, A_{t}) + \sum_{k=t}^{min(t+n-1, T-1)} \delta_{k}\prod_{i=t+1}^{k} \gamma\pi(A_{i}|S_{i})$.
\\

With some clever math and using 7.8, 7.16 can be rewritten as: 
\\

$G_{t:t+n} = R_{t+1} + \gamma(\bar{V}(S_{t+1}) + \pi(A_{t+1}|S_{t+1})(G_{t+1} - Q(S_{t+1}, A_{t+1})))$
\\

and then subtract $Q(S_{t}, A_{t})$ to get: 
\\

$G_{t:t+n} - Q(S_{t}, A_{t}) = R_{t+1} + \gamma\bar{V}(S_{t+1}) + \pi(A_{t+1}|S_{t+1})(G_{t+1} - 

Q(S_{t+1}, A_{t+1}))) - Q(S_{t}, A_{t})$
\\

$= (R_{t+1} + \gamma\bar{V}(S_{t+1}) - Q(S_{t}, A_{t})) + \gamma \pi(A_{t+1}|S_{t+1})(G_{t+1} - 

Q(S_{t+1}, A_{t+1}))$
\\

$= \bar{\delta_{t}} + \gamma \pi(A_{t+1}|S_{t+1})(G_{t+1} - Q(S_{t+1}, A_{t+1}))$\\

Now take into account the term $(G_{t+1} - Q(S_{t+1}, A_{t+1}))$ -- we can expand as 

before:
\\

$G_{t+1} - Q(S_{t+1}, A_{t+1}) = \bar{\delta}_{t+1} + \gamma \pi(A_{t+2}|S_{t+2})(G_{t+2} - Q(S_{t+2}, A_{t+2}))$
\\

and substitute back into the original express to get:
\\

$G_{t:t+n} - Q(S_{t}, A_{t}) = \bar{\delta}_{t} + \gamma \pi(A_{t+1}|S_{t+1})(\bar{\delta}_{t+1} + \gamma \pi(A_{t+2}|S_{t+2})(G_{t+2} - Q(S_{t+2}, A_{t+2})))
\\

= $\bar{\delta}_{t} + \gamma \pi(A_{t+1}|S_{t+1})(\bar{\delta}_{t+1}) + \gamma^{2}(\pi(A_{t+1}|S_{t+1})\pi(A_{t+2}|S_{t+2}))(G_{t+2} - Q(S_{t+2}, A_{t+2})).$
\\ 

this ultimately yields the original equation:
\\

$G_{t:t+n} = Q(S_{t}, A_{t}) + \sum_{k=t}^{min(t+n-1, T-1)} \delta_{k}\prod_{i=t+1}^{k} \gamma\pi(A_{i}|S_{i}).$









\end{document}
